\documentclass{article}

%% Inclusion of images
\usepackage{graphicx}

%% Authors names and institutions
\usepackage{authblk}

%% Drawing
\usepackage{tikz}
\usetikzlibrary{arrows,calc}

% Tikz style for graph drawing
\tikzstyle{arc}=[->,shorten <=3pt, shorten >=3pt,
                 >=stealth, line width=1.1pt]
\tikzstyle{edge}=[shorten <=2pt, shorten >=2pt,
                  >=stealth, line width=1.1pt]
\tikzstyle{vertex}=[circle, fill=white, draw,
                    minimum size=5pt,
                    inner sep=0pt, outer sep=0pt]

%% Positioning of floats
\usepackage{float}

%% Color in text, tables, etc.
\usepackage{xcolor}

%% Additional symbols
\usepackage{amsmath}
\usepackage{amssymb}

%% Theorem environments
\usepackage{amsthm}

%% References with automatic names
\usepackage[capitalize]{cleveref}

% Theorem definition
\newtheorem{theorem}{Theorem}
\newtheorem{lemma}[theorem]{Lemma}
\newtheorem{corollary}[theorem]{Corollary}
\newtheorem{proposition}[theorem]{Proposition}

\title{Title%
\thanks{The authors gratefully acknowledge support from some grant}}

\author[1]{First~Author\thanks{first@mail.edu}}
\author[2]{Second~Author\thanks{second@mail.mx}}
\author[3]{Third~Author\thanks{third@mail.go}}

\affil[1]{First Author's Department,
          First Author's Univeristy}
\affil[2]{Second Author's Department,
          Second Author's Univeristy}
\affil[3]{Third Author's Department,
          Third Author's Univeristy}

\begin{document}
\date{}

\maketitle
\begin{abstract}
    Lorem ipsum dolor sit amet, consectetur adipisicing elit, sed do eiusmod
    tempor incididunt ut labore et dolore magna aliqua. Ut enim ad minim veniam,
    quis nostrud exercitation ullamco laboris nisi ut aliquip ex ea commodo
    consequat. Duis aute irure dolor in reprehenderit in voluptate velit esse
    cillum dolore eu fugiat nulla pariatur. Excepteur sint occaecat cupidatat
    non proident, sunt in culpa qui officia deserunt mollit anim id est laborum.
\end{abstract}


\section{Introduction}

Try using the \texttt{cleveref} package :D

\begin{theorem}
\label{thm:first}
First example theorem.
\end{theorem}

\begin{theorem}
\label{thm:second}
Second example theorem.
\end{theorem}

\begin{lemma}
\label{lem:first}
First example lemma.
\end{lemma}

\texttt{Cleveref} automatically detects the type of environment you
are referring to, and fills in its name.   For example, you might
refer to \Cref{thm:first,thm:second}, or to \Cref{thm:second,lem:first}.


\vspace{2mm}


\begin{thebibliography}{}

\bibitem{bang2009}
  J.~Bang-Jensen and G.~Gutin,
  Digraphs Theory Algorithms and Applications
  Second Edition,
  Springer-Verlag London (2009).

\bibitem{haskinsSIAMJC2}
  L. Haskins, and D.J. Rose,
  Toward characterization of perfect elimination digraphs,
  SIAM Journal on Computing 2 (1973) 217 -- 224.
\end{thebibliography}


\end{document}
